\addtocontents{toc}{\protect\Needspace{3\baselineskip}}
\originalchapter{扰动增长与结构形成}\label{ch:structure}
\originalsection{线性扰动}
写 $\rho=\bar\rho(1+\delta)$, 把物理速度分为 Hubble 流与特殊速度 $\boldsymbol u$. 在共动坐标中连续、Euler 与 Poisson 方程线性化为
\begin{align*}
\dot\delta+\frac1a\nabla\cdot\boldsymbol u&=0,\\
\dot{\boldsymbol u}+H\boldsymbol u&=-\frac1a\nabla\phi-\frac{c_s^2}a\nabla\delta,\\
\nabla^2\phi&=4\pi Ga^2\bar\rho\delta.
\end{align*}
对第一式求导, 消去速度与势, Fourier 模满足
\begin{equation}\label{eq:growth}
\ddot\delta_{\boldsymbol k}+2H\dot\delta_{\boldsymbol k}
+\left(\frac{c_s^2k^2}{a^2}-4\pi G\bar\rho\right)\delta_{\boldsymbol k}=0.
\end{equation}
此 Newton 近似用于远小于视界、非相对论组分的尺度, 超视界扰动需要相对论且必须指定规范. 在 Einstein--de Sitter 尘埃模型中 $H=2/(3t)$、$4\pi G\bar\rho=2/(3t^2)$, 代 $\delta\propto t^p$ 得 $p=2/3,-1$, 生长模 $\delta\propto a$. 辐射主导时物质扰动的增长受背景抑制, 暗能量主导后也不再近似 $D=a$.

\originalsection{统计描述}
取 $\delta(\boldsymbol x)=\int\dd^3k\,\delta(\boldsymbol k)\ee^{\ii\boldsymbol k\cdot\boldsymbol x}/(2\pi)^3$, 齐次场满足 $\langle\delta(\boldsymbol k)\delta^*(\boldsymbol k')\rangle=(2\pi)^3\delta_D(\boldsymbol k-\boldsymbol k')P(k)$. 相关函数 $\xi(r)$ 是 $P$ 的 Fourier 变换; $\Delta^2(k)=k^3P(k)/(2\pi^2)$ 是每对数波数的方差贡献. 原初谱经转移函数与生长因子传播, $P(k,t)=D(t)^2T(k)^2P_{\rm initial}(k)$ 只在适用的线性范围成立.

给定平滑窗口 $W(kR)$, $\sigma^2(R)=\int\Delta^2(k)|W(kR)|^2\dd\ln k$. 对实空间球顶帽 $W(x)=3(\sin x-x\cos x)/x^3$, 小 $x$ 极限为 1. 窗口及质量映射 $M=4\pi\bar\rho_0R^3/3$ 必须一起给出, 否则不同 $\sigma(M)$ 不能直接比较.

\originalsection{球对称坍缩}
均匀球顶帽在无压尘埃、无壳交叉条件下, 每层满足 $\ddot r=-GM/r^2$. 负能解写成 $r=A(1-\cos\eta)$、$t=B(\eta-\sin\eta)$、$A^3=GMB^2$. 在 $\eta=\pi$ 达最大半径, 在 $2\pi$ 形式上坍为一点. EdS 背景密度为 $1/(6\pi Gt^2)$, 代入球平均密度得到转向密度比 $9\pi^2/16$. 早期展开给线性外推 $\delta_L=(3/5)[3(\eta-\sin\eta)/4]^{2/3}$, 在形式坍缩时为 $\delta_c\simeq1.686$.

若能量守恒且最终满足维里关系, $W=-3GM^2/(5R)$ 导致 $R_{\rm vir}=R_{\rm ta}/2$. 因背景密度在 $t_{\rm coll}=2t_{\rm ta}$ 时下降四倍, 晕密度因半径减半提高八倍, 最终密度比为 $18\pi^2$. 这些数值是 EdS 理想化, 不可无条件用于任意红移的 $\Lambda$CDM.

\originalsection{质量函数与冷却}
Gaussian 平滑场超过阈值的体积分数是 $F(>M)=\frac12\operatorname{erfc}[\delta_c/(\sqrt2\sigma(M))]$. 在 $\sigma\to\infty$ 的小质量极限只到 $1/2$, 原始 Press--Schechter 计数漏掉嵌套区域. 锐波数滤波使 $\delta$ 随方差 $S$ 作 Markov 随机游走. 将阈值设为吸收壁, 镜像法从自由 Gaussian 核减去镜像核, 首次过壁分布为 $f(S)=\delta_c\exp[-\delta_c^2/(2S)]/(\sqrt{2\pi}S^{3/2})$, 积分为 1. 这解释因子 2, 但实空间顶帽游走并非 Markov, 椭球坍缩和数值晕质量函数还需更多修正.

气体能否形成恒星还取决于冷却. 光学薄气体损失率写作 $n_en_i\Lambda(T)$, 储能约为 $3n\kb T/2$, 因而 $t_{\rm cool}=3n\kb T/(2n_en_i\Lambda)$. 与 $t_{\rm dyn}\sim(G\rho)^{-1/2}$ 比较可估计哪些晕能快速形成冷气体. $\Lambda$ 的单位和电离、组成约定不能省略, 固定一常数冷却曲线只是数量级模型.

\originalsection{坍缩阈值的早期展开}
球顶帽的密度比可以直接用参数解而不用背诵常数. $\rho_{\rm halo}=3M/[4\pi A^3(1-\cos\eta)^3]$, EdS 背景 $\bar\rho=1/(6\pi GB^2(\eta-\sin\eta)^2)$, 结合 $A^3=GMB^2$ 得
\[
1+\delta=\frac92\frac{(\eta-\sin\eta)^2}{(1-\cos\eta)^3}.
\]
展开 $\eta-\sin\eta=\eta^3/6-\eta^5/120+\cdots$、$1-\cos\eta=\eta^2/2-\eta^4/24+\cdots$, 比值给 $\delta=3\eta^2/20+O(\eta^4)$. 同时 $t\simeq B\eta^3/6$, 线性增长 $t^{2/3}$ 因而唯一匹配为 $\delta_L=(3/20)[6(\eta-\sin\eta)]^{2/3}$, 与正文写法等价. 在 $\eta=2\pi$ 外推得 $\delta_c=(3/20)(12\pi)^{2/3}=1.68647$.

外推的 $\delta_L$ 不是坍缩时真实密度对比, 后者在理想无压模型发散. 两者用途不同: 前者把线性随机场与形成条件联系, 后者描述非线性形状. Press--Schechter 使用的是前者. 后续维里化并非由 $\eta=2\pi$ 的点状解自行产生, 而是额外加入能量守恒和稳态维里假设后用有限半径模型替换.
